% ECONOMICS_PROBLEM: {"id": "D44-83", "title": "Symmetric equilibrium existence with multidimensional auction information", "jel_code": "D44", "source_id": "OP-0592"}
% FIRST_POSTED: 2026-10-09
% ATTEMPT_STATUS: unresolved
% ENTRY_KIND: research_attempt
\documentclass[11pt]{article}
\usepackage[T1]{fontenc}
\usepackage[utf8]{inputenc}
\usepackage[margin=27mm]{geometry}
\usepackage{amsmath,amssymb,amsthm}
\usepackage[hidelinks]{hyperref}
\newtheorem{theorem}{Theorem}
\newtheorem{proposition}[theorem]{Proposition}
\newtheorem{lemma}[theorem]{Lemma}
\newtheorem{corollary}[theorem]{Corollary}
\theoremstyle{definition}
\newtheorem{definition}[theorem]{Definition}
\newtheorem{example}[theorem]{Example}
\theoremstyle{remark}
\newtheorem{remark}[theorem]{Remark}
\setlength{\parskip}{4pt}
\title{Multidimensional auction information: absence of bid atoms and global deviation certificates}
\author{GPT 6 Astra Ultra}
\date{9 October 2026}
\begin{document}
\maketitle
\noindent\textbf{Problem:} \texttt{D44-83}. \textbf{Status:} Partial results; the complete catalogue target remains unresolved.
\section{Problem and scope}
The remaining catalogue question allows jumps and pooling and asks for symmetric pure equilibrium under the prescribed auction and tie rule. The previous GPT 6 Astra Pro note treats a binary-signal family and its continuous-strategy threshold. This note does not repeat that classification. It proves that pooling is impossible in a positive-value two-bidder subclass, and gives a rigorous way to convert finite deviation checks into a global bound. Neither result proves general existence.

\section{Two-bidder primitives}
There is one object and two symmetric bidders. Quality $q\in[0,1]$ has positive density $f$, and independent tastes $t_i\in[0,\bar t]$ have positive continuous density $w$. Taste and quality are independent. Signals are conditionally independent with probabilities $\alpha_s(q)\beta_s(t)$, which sum to one for every $(q,t)$. The finite family $\alpha_s$ is positive and bounded away from zero, and $\beta_s$ is nonnegative measurable. Bidders value the object at bounded continuous $u(q,t)$, strictly increasing in $t$ for each $q$. For this note impose the additional restriction
\[
 u(q,t)>0\quad\text{for Lebesgue-almost every }q\text{ and every supported taste }t.
\]
The highest bid wins and pays the other bid; a tie is split with probability $1/2$. All bids are nonnegative and finite. Write $b_s(t)$ for a measurable symmetric pure strategy; monotonicity is not needed in the next theorem.

\begin{lemma}[A tie posterior does not depend on own taste]
Suppose the opponent's bid has an unconditional atom at $z$. Put
\[
 c_r(z)=\int_{\{t:b_r(t)=z\}}w(t)\beta_r(t)\,dt,
 \qquad A_z(q)=\sum_r\alpha_r(q)c_r(z).
\]
Then $A_z(q)>0$ for every $q$. Given own supported signal $s$, taste $t$, and opponent bid $z$, the posterior density of quality is proportional to $f(q)\alpha_s(q)A_z(q)$, independently of $t$.
\end{lemma}
\begin{proof}
An unconditional atom implies $c_r(z)>0$ for at least one $r$, and positivity of $\alpha_r$ gives the first assertion. Conditional independence gives joint density proportional to
\[
 f(q)\alpha_s(q)\beta_s(t)w(t)A_z(q).
\]
Conditioning on a supported own $(s,t)$ cancels the factors involving $t$, leaving the asserted density. The normalizing integral is positive.
\end{proof}

\begin{theorem}[No atoms in a symmetric pure equilibrium]
Under these primitives, the unconditional bid distribution of any symmetric pure Bayesian Nash equilibrium is atomless. Therefore no signal strategy can pool a set of tastes of positive conditional probability at one bid. Weakly increasing equilibrium strategies may still have jumps.
\end{theorem}
\begin{proof}
Suppose an atom exists at $z>0$. For almost every type $(t,s)$ bidding $z$, consider bids $z+\varepsilon$ and $z-\varepsilon$, with $0<\varepsilon<z$. As $\varepsilon\downarrow0$, the payoff changes from opponents bidding strictly between the old and new bids vanish by dominated convergence. On an exact tie, upward deviation changes the winning probability from $1/2$ to $1$, and downward deviation changes it from $1/2$ to zero. Best-response inequalities therefore imply both signs of the same conditional surplus, so
\[
 \mathbb E[u(q,t)\mid s,B_{-i}=z]=z.
\]
By the lemma the conditional law in this equation is independent of $t$. Its left-hand side is strictly increasing in $t$: for $t'<t$, the integrand is strictly larger at every $q$. The equality can hold for at most one taste for each signal. But a finite union of such singleton tastes has zero probability under densities $w(t)\beta_s(t)$, contradicting the positive atom.

If $z=0$, only the upward deviation is needed. On a tie its limiting gain is half the positive conditional expectation of $u(q,t)$, multiplied by the positive tie probability. The strict positivity assumption ensures that gain is positive, again contradicting optimality. Thus no nonnegative atom exists. A jump of a strictly increasing function omits a bid interval without concentrating positive taste probability, so the result does not exclude jumps.
\end{proof}

\section{Finite checks with a continuum guarantee}
Now take any measurable proposed symmetric profile, not necessarily an equilibrium. Assume its conditional opponent bid distributions are atomless. For every supported signal $s$ they are independent of own taste, by the same factorization argument. Suppose $0\le u\le U$ and the proposed bids lie in $[0,B]$ with $B\ge U$. Let $G$ be a finite grid in $[0,B]$ containing both endpoints and with largest spacing $h$. Define
\[
 \omega_s(h)=\sup_x\Pr(B_{-i}\in[x,x+h]\mid s).
\]

\begin{proposition}[Global incentive certificate]
If for almost every $(t,s)$ the profile satisfies
\[
 \Pi(t,s;b_s(t))\ge \max_{g\in G}\Pi(t,s;g)-\varepsilon,
\]
then its gain from any nonnegative bid deviation is at most
\[
 \varepsilon+B\max_s\omega_s(h).
\]
If conditional bid densities exist and are bounded by $L_s$, this is at most $\varepsilon+Bh\max_s L_s$.
\end{proposition}
\begin{proof}
Against an atomless opponent, bids $x,y\in[0,B]$ differ in payoff only when the opponent bid lies between them. The absolute winning surplus on this event is at most $B$, since both the value and opponent bid are in $[0,B]$. Thus $|\Pi(t,s;x)-\Pi(t,s;y)|\le B\omega_s(h)$ whenever $|x-y|\le h$. Approximate any deviation in $[0,B]$ by a grid point within $h$ and use the assumed inequality. A bid greater than $B$ changes winning only against opponent bids at or above $B$, an event of probability zero here; alternatively, capping at $B\ge U$ never loses positive surplus. The density bound follows by integration over an interval of length $h$.
\end{proof}

\section{Remaining gap and use of the certificate}
The positivity restriction matters: with negative values and a nonnegative bid floor, zero-bid pooling cannot be ruled out by the upward-deviation argument. For more than two bidders, upward and downward tie deviations weight different pivotal tie configurations, so the scalar equality used above needs a new argument. Even in the positive two-bidder case, atomlessness supplies no compactness or fixed-point theorem in the monotone pure-strategy class. Finally, grid checks must hold for the continuum of own types, or have a separate certified interpolation bound. A finite sample of own types and a finite sample of deviations alone are not a proof. Atomic discretizations require explicit one-sided tie checks and do not satisfy the proposition's assumptions.

\begin{thebibliography}{9}
\bibitem{PS} W. Pesendorfer and J. M. Swinkels (2000), Efficiency and information aggregation in auctions, \emph{American Economic Review} 90(3), 499--525. \url{https://www.aeaweb.org/articles?id=10.1257/aer.90.3.499}. Source of the multidimensional information framework; no general finite-auction pure-existence conclusion is inferred here.
\bibitem{ACM} A. Araujo, L. I. de Castro Filho and H. Moreira (2004), \emph{Pure Strategy Equilibria of Multidimensional and Non-Monotonic Auctions}, working paper. \url{https://epge.fgv.br/files/1608.pdf}. Related characterization and tie-rule results; different type-independence and tie assumptions must be checked before importing its existence claims.
\bibitem{prior} GPT 6 Astra Pro (2026), \emph{Continuous-Strategy Nonexistence in a Two-Bidder Auction}, companion research attempt for D44-83, dated 8 October 2026. \url{https://github.com/mehmetmars7/Erdosproblems-llm-hunter/blob/main/attacks/open_problems/economics/gpt_6_astra_pro/D44-83.tex}. A prior generated research note, not a peer-reviewed source; its threshold classification is not used in either proof above.
\end{thebibliography}
\end{document}
